% Lettre Universite Laval - meme modele typographique que la lettre HEC

\documentclass[12pt,a4paper]{article}

\usepackage[a4paper,top=1.9cm,bottom=1.9cm,left=2.4cm,right=2.4cm]{geometry}
\usepackage[utf8]{inputenc}
\usepackage[T1]{fontenc}
\usepackage[scaled=1.0]{helvet}
\renewcommand{\familydefault}{\sfdefault}
\usepackage{xcolor}
\usepackage[hidelinks]{hyperref}
\usepackage{titlesec}

\definecolor{accent}{HTML}{1F4E79}
\definecolor{muted}{HTML}{5A6165}
\definecolor{rule}{HTML}{D5E2EC}

\linespread{1.5}
\setlength{\parindent}{0pt}
\setlength{\parskip}{0.3em}
\pagestyle{plain}

\titleformat{\section}
  {\normalsize\bfseries\color{accent}}
  {}{0pt}{}
\titlespacing*{\section}{0pt}{0.9em}{0.15em}

\begin{document}
\thispagestyle{empty}

{\LARGE\bfseries Stéphane KPOVIESSI}\par
\vspace{0.25em}
{\footnotesize\color{muted}Cergy / Île-de-France \textbullet{} +33 7 48 59 47 13 \textbullet{} oastephanekpoviessi@gmail.com\\
linkedin.com/in/stephanekpoviessi \textbullet{} seshat-ai.com}\par
\vspace{0.25em}
{\color{rule}\hrule height 1pt}
\vspace{0.3em}

{\bfseries Lettre de motivation \textbullet{} MBA -- Analytique d'affaires\\
Université Laval \textbullet{} Automne 2027}\par
\vspace{0.5em}

Madame, Monsieur,

Je présente ma candidature au MBA en analytique d'affaires de l'Université Laval pour l'automne 2027. Ingénieur diplômé de JUNIA ISEN, spécialisé en Big Data et Data Science, je souhaite compléter mes compétences techniques par une formation en administration des affaires. Mon objectif est de relier les analyses aux choix stratégiques et financiers, puis de conduire des projets d'analytique au service de la performance.

\section{Un parcours scientifique tourné vers la décision}

J'ai suivi un seul cursus d'ingénieur de cinq années consécutives, de 2020 à 2025 : deux années préparatoires à Sainte Marie-Stella au Bénin, une troisième année à JUNIA ISEN Rabat, puis deux années à JUNIA ISEN Lille. L'attestation délivrée en avril 2026 fait référence au jury de septembre 2025. Cette formation m'a donné des bases quantitatives solides : en deuxième année préparatoire, ma moyenne du second semestre était de 16,94/20, avec 19,70/20 en mathématiques et 19,45/20 en informatique. J'ai ensuite approfondi l'aide à la décision, la visualisation, l'apprentissage automatique et le traitement de données à grande échelle.

Chez Allianz France, j'ai participé à la migration d'un datamart d'assurance de SAS vers Python/PySpark et à la fiabilisation d'indicateurs métier. Chez Lyne \& Frères, mes analyses commerciales, opérationnelles et agricoles ainsi que mes tableaux de bord ont contribué au suivi de l'activité. Mon expérience de R\&D en vision par ordinateur chez Sylvagreg a renforcé ma pratique de l'expérimentation. Ces expériences m'ont appris à adapter les analyses aux besoins des équipes. Elles m'ont aussi donné envie de mieux comprendre les arbitrages de gestion auxquels elles contribuent.

\newpage
\section{Pourquoi le MBA en analytique d'affaires à Laval}

Je choisis l'Université Laval pour l'articulation entre formation générale en gestion et exploitation des données. Les cours Information financière et prise de décision, Finance corporative et Stratégie et structure de l'entreprise répondent aux compétences que je souhaite développer. Je veux apprendre à évaluer un projet en tenant compte de ses objectifs, de ses ressources et de ses conséquences financières.

Les enseignements Analytique d'affaires et Systèmes d'aide à la décision prolongent directement mon parcours. L'activité terminale, sous forme de projet ou de stage, m'intéresse pour confronter ces acquis à une situation réelle. J'aimerais notamment travailler sur des indicateurs reliant performance commerciale, opérations et priorités de développement, en tenant compte de la qualité des données disponibles.

\section{Un projet professionnel et entrepreneurial cohérent}

À court terme, je vise un poste d'analyste en intelligence d'affaires ou de conseiller en analytique et en performance. Je souhaite prendre en charge le cadrage des besoins et formuler des recommandations utilisables par les décideurs, puis coordonner des projets de transformation. Mon projet entrepreneurial Seshat, en développement depuis septembre 2026, renforce cette ambition : concevoir des solutions d'IA et d'automatisation suppose aussi de comprendre les usages, de structurer une offre et d'arbitrer des ressources. À long terme, je souhaite consolider ce projet et contribuer au développement de Lyne \& Frères.

Mes responsabilités de responsable communication de JUNIA ALUMNI en 2023--2024, puis de vice-président en 2024--2025, m'ont initié à la coordination et au travail collectif. Je souhaite enrichir les échanges à Laval par cette expérience, mon parcours dans trois pays et ma perspective d'ingénieur.

Je vous remercie de l'attention portée à ma candidature et me tiens à votre disposition pour tout complément d'information. Veuillez agréer, Madame, Monsieur, l'expression de mes salutations distinguées.

\vspace{0.4em}
{\bfseries Stéphane KPOVIESSI}
\end{document}
